\subsection{R 计算：边缘分布}

边缘分布在 R 中通过逐行或逐列求和（离散）或积分（连续）计算。

\begin{example}[离散边缘分布]\label{ex:marginal-discrete-ch2}
设 $(X_1, X_2)$ 的联合 pmf 如下（列联表形式）：

\[
\begin{array}{c|cccc|c}
X_2 \backslash X_1 & 1 & 2 & 3 & 4 & \text{边缘} \\ \hline
1 & 0.05 & 0.10 & 0.05 & 0.02 & ? \\
2 & 0.08 & 0.15 & 0.12 & 0.05 & ? \\
3 & 0.02 & 0.08 & 0.10 & 0.08 & ? \\ \hline
\text{边缘} & ? & ? & ? & ? & 1.00
\end{array}
\]

则 $X_1$ 的边缘 pmf 为每列之和：
\[
p_{X_1}(1) = 0.05 + 0.08 + 0.02 = 0.15, \quad
p_{X_1}(2) = 0.10 + 0.15 + 0.08 = 0.33, \quad \text{等等}
\]

R 计算：
\begin{verbatim}
# 联合 pmf（矩阵形式）
joint <- matrix(c(0.05, 0.08, 0.02,  # X1=1
                  0.10, 0.15, 0.08,  # X1=2
                  0.05, 0.12, 0.10,  # X1=3
                  0.02, 0.05, 0.08), # X1=4
                nrow = 3, byrow = FALSE)
colnames(joint) <- c("X1=1", "X1=2", "X1=3", "X1=4")
rownames(joint) <- c("X2=1", "X2=2", "X2=3")

# 边缘分布（按行、按列求和）
margin_x1 <- colSums(joint)  # X1 的边缘 pmf
margin_x2 <- rowSums(joint)  # X2 的边缘 pmf

# 验证：边缘分布之和为 1
sum(margin_x1)  # = 1
sum(margin_x2)  # = 1
\end{verbatim}
\end{example}

\begin{example}[连续边缘分布]\label{ex:marginal-continuous-ch2}
设 $(X_1, X_2)$ 的联合 pdf 为
\[
f(x_1, x_2) = \begin{cases}
8x_1 x_2 & 0 < x_1 < x_2 < 1, \\
0 & \text{otherwise}.
\end{cases}
\]

求 $X_1$ 的边缘 pdf。

由于 $0 < x_1 < x_2 < 1$，$x_2$ 的取值范围是 $x_1 < x_2 < 1$，因此：
\[
f_{X_1}(x_1) = \int_{x_1}^{1} 8x_1 x_2 \, dx_2 = 8x_1 \left[\frac{x_2^2}{2}\right]_{x_2=x_1}^{x_2=1} = 8x_1 \cdot \frac{1 - x_1^2}{2} = 4x_1(1 - x_1^2),
\]
其中 $0 < x_1 < 1$。

验证：
\[
\int_0^1 4x_1(1 - x_1^2) \, dx_1 = 4\left[\frac{x_1^2}{2} - \frac{x_1^4}{4}\right]_0^1 = 4\left(\frac{1}{2} - \frac{1}{4}\right) = 1.
\]

R 数值验证：
\begin{verbatim}
# 验证边缘 pdf 的数值积分
f_marginal_x1 <- function(x1) {
  f <- function(x2) 8 * x1 * x2
  integrate(f, lower = x1, upper = 1)$value
}

# 在几个点验证
x1_vals <- c(0.2, 0.5, 0.8)
sapply(x1_vals, f_marginal_x1)  # 应该接近 4*x1*(1-x1^2)
4 * x1_vals * (1 - x1_vals^2)
\end{verbatim}
\end{example}

\subsection{边缘分布的深度理解}

\textbf{为什么边缘分布重要？} 边缘分布在实际应用中有广泛用途：

\textbf{1. 人口统计学}：当我们研究"收入"与"教育年限"的联合分布时，边缘分布给出了各变量的边缘分布——即不考虑另一变量时该变量的分布。这相当于"总体"的分布。

\textbf{2. 保险精算}：在汽车保险中，我们关心"理赔次数"与"理赔金额"的联合分布，但边缘分布分别给出了理赔次数的分布和理赔金额的分布——这两个分布本身就有独立的应用价值。

\textbf{3. 信号处理}：在处理二元信号时，边缘分布给出了各分量信号的分布，这在单独分析各分量时有用。

\textbf{信息损失}：一个关键事实是，从联合分布到边缘分布会丢失信息。如果只知道边缘分布，则无法恢复联合分布——因为这需要知道变量之间的依赖结构。

\textbf{Clayton 模型}：在金融保险中，Clayton copula 正是用来说明：相同的边缘分布可以对应于不同的联合分布。这个事实在风险管理中非常重要——仅知道各变量的边缘分布是不够的。
